%==============================================================================
\section{Khái niệm nền --- Thống kê cho học máy}
\label{sec:khai-niem-thong-ke}
%==============================================================================

\begin{dinhnghia}[title={Thống kê (Statistics)}]
Thống kê là nhánh toán học về thu thập, phân tích, diễn giải và trình bày dữ liệu.
Nó cung cấp phương pháp và kỹ thuật để phân tích dữ liệu và rút kết luận.
\end{dinhnghia}

\begin{ynghia}[title={Vai trò trong ML}]
Thống kê là nền tảng của học máy: phân tích và trực quan hoá để tìm pattern ẩn.
Dùng trong kiểm định mô hình, làm sạch dữ liệu, chọn mô hình, đánh giá hiệu năng, \ldots
\end{ynghia}

\begin{phuongphap}[title={Các khái niệm trong phần này}]
Mean, median, mode; độ lệch chuẩn, phương sai; percentiles; phân phối dữ liệu; skewness và kurtosis; bias và variance; giả thuyết (hypothesis).
\end{phuongphap}

\begin{phuongphap}[title={Hai nhánh thống kê}]
\begin{itemize}
  \item \textbf{Thống kê mô tả (descriptive)}: tóm tắt đặc trưng tập dữ liệu (mean, std, \ldots).
  \item \textbf{Thống kê suy luận (inferential)}: suy luận về quần thể từ mẫu (kiểm định, hồi quy, khoảng tin cậy).
\end{itemize}
\end{phuongphap}
